\documentclass[aps,prb,twocolumn,superscriptaddress]{revtex4-2}
\usepackage{xeCJK}
\setCJKmainfont{SimSun}[ItalicFont=KaiTi, BoldFont=SimHei]
\setCJKsansfont{SimHei}
\setCJKmonofont{SimSun}
\usepackage{algorithm}
\usepackage{algorithmic}
\renewcommand{\algorithmicrequire}{\textbf{输入:}}
\renewcommand{\algorithmicensure}{\textbf{输出:}}
\floatname{algorithm}{算法}
\usepackage{graphicx}
%\usepackage{subcaption}
%\usepackage[draft]{graphicx}
\usepackage{latexsym}
\usepackage{amssymb}
\usepackage{amsmath}
\usepackage{amsfonts}
\usepackage[vcentermath]{youngtab}
\usepackage{upgreek}
\usepackage{bm,dsfont}
\usepackage{multirow}
\usepackage{enumitem}
\usepackage{color}
\usepackage{hyperref}
\hypersetup{
colorlinks = true,
linkcolor = [rgb]{0.70,0.13,0.13},
citecolor = [rgb]{0.13,0.55,0.13},
urlcolor = [rgb]{0.25, 0.41, 0.88}}
\usepackage{makecell}
\renewcommand\theadalign{bc}
\renewcommand\theadfont{\bfseries}
\renewcommand\theadgape{\Gape[0pt]}
\renewcommand\cellgape{\Gape[2pt]}
\newcommand{\bra}[1]{\langle #1|}
\newcommand{\ket}[1]{|#1 \rangle}
\newcommand{\braket}[2]{\langle #1|#2 \rangle}
\newcommand{\dd}{\mathrm{d}}
\newcommand{\DD}{\mathcal{D}}
\newcommand{\ii}{\mathrm{i}}
\newcommand{\e}{\mathrm{e}}
\newcommand{\rep}{\vect{r}}
\newcommand{\ybox}{\tiny\yng}
\newcommand{\id}{\mathds{1}}
\newcommand{\E}{\mathop{\mathbb{E}}}
\newcommand{\LU}{\mathrm{LU}}
\newcommand{\U}{\mathrm{U}}
\newcommand{\SU}{\mathrm{SU}}
\renewcommand{\O}{\mathrm{O}}
\newcommand{\SO}{\mathrm{SO}}
\newcommand{\Sp}{\mathrm{Sp}}
\newcommand{\Spin}{\mathrm{Spin}}
\renewcommand{\u}{\mathfrak{u}}
\newcommand{\su}{\mathfrak{su}}
\renewcommand{\o}{\mathfrak{o}}
\newcommand{\so}{\mathfrak{so}}
\newcommand{\spin}{\mathfrak{spin}}
\newcommand{\Cl}{{\mathcal{C}\ell}}
\newcommand{\dsA}{\mathbb{A}}
\newcommand{\dsB}{\mathbb{B}}
\newcommand{\dsC}{\mathbb{C}}
\newcommand{\dsD}{\mathbb{D}}
\newcommand{\dsE}{\mathbb{E}}
\newcommand{\dsF}{\mathbb{F}}
\newcommand{\dsG}{\mathbb{G}}
\newcommand{\dsH}{\mathbb{H}}
\newcommand{\dsI}{\mathbb{I}}
\newcommand{\dsJ}{\mathbb{J}}
\newcommand{\dsK}{\mathbb{K}}
\newcommand{\dsL}{\mathbb{L}}
\newcommand{\dsM}{\mathbb{M}}
\newcommand{\dsN}{\mathbb{N}}
\newcommand{\dsO}{\mathbb{O}}
\newcommand{\dsP}{\mathbb{P}}
\newcommand{\dsQ}{\mathbb{Q}}
\newcommand{\dsR}{\mathbb{R}}
\newcommand{\dsS}{\mathbb{S}}
\newcommand{\dsT}{\mathbb{T}}
\newcommand{\dsU}{\mathbb{U}}
\newcommand{\dsV}{\mathbb{V}}
\newcommand{\dsW}{\mathbb{W}}
\newcommand{\dsX}{\mathbb{X}}
\newcommand{\dsY}{\mathbb{Y}}
\newcommand{\dsZ}{\mathbb{Z}}
\newcommand{\scA}{\mathcal{A}}
\newcommand{\scB}{\mathcal{B}}
\newcommand{\scC}{\mathcal{C}}
\newcommand{\scD}{\mathcal{D}}
\newcommand{\scE}{\mathcal{E}}
\newcommand{\scF}{\mathcal{F}}
\newcommand{\scG}{\mathcal{G}}
\newcommand{\scH}{\mathcal{H}}
\newcommand{\scI}{\mathcal{I}}
\newcommand{\scJ}{\mathcal{J}}
\newcommand{\scK}{\mathcal{K}}
\newcommand{\scL}{\mathcal{L}}
\newcommand{\scM}{\mathcal{M}}
\newcommand{\scN}{\mathcal{N}}
\newcommand{\scO}{\mathcal{O}}
\newcommand{\scP}{\mathcal{P}}
\newcommand{\scQ}{\mathcal{Q}}
\newcommand{\scR}{\mathcal{R}}
\newcommand{\scS}{\mathcal{S}}
\newcommand{\scT}{\mathcal{T}}
\newcommand{\scU}{\mathcal{U}}
\newcommand{\scV}{\mathcal{V}}
\newcommand{\scW}{\mathcal{W}}
\newcommand{\scX}{\mathcal{X}}
\newcommand{\scY}{\mathcal{Y}}
\newcommand{\scZ}{\mathcal{Z}}
\newcommand{\Tr}{\operatorname{Tr}}
\newcommand{\sgn}{\operatorname{sgn}}
\newcommand{\vol}{\operatorname{vol}}
\renewcommand{\Re}{\operatorname{Re}}
\renewcommand{\Im}{\operatorname{Im}}
\newcommand{\vect}[1]{{\bm{#1}}}
\newcommand{\norm}[1]{{\lVert #1\rVert}}
\newcommand{\mat}[1]{\left[\begin{matrix}#1\end{matrix}\right]}
\newcommand{\smat}[1]{\left[\begin{smallmatrix}#1\end{smallmatrix}\right]}
\newcommand{\eq}[1]{\begin{equation}#1\end{equation}}
\newcommand{\eqs}[1]{\begin{equation}\begin{split}#1\end{split}\end{equation}}
\newcommand{\eqnref}[1]{式\,\eqref{#1}}
\newcommand{\figref}[1]{图\,\ref{#1}}
\newcommand{\tabref}[1]{表\,\ref{#1}}
\newcommand{\secref}[1]{第\,\ref{#1}\,节}
\newcommand{\appref}[1]{附录\,\ref{#1}}
\newcommand{\refcite}[1]{文献\,\onlinecite{#1}}
\newcommand{\tocite}[1]{[{\color[rgb]{0.13,0.55,0.13}{#1}}]}

\def \Z{\mathbb{Z}}

\newcommand{\bea}{\begin{eqnarray}}
\newcommand{\eea}{\end{eqnarray}}
\def\be{\begin{equation}}
\def\ee{\end{equation}}
\newcommand{\beq}{\begin{equation}}
\newcommand{\eeq}{\end{equation}}
\newcommand{\beqn}{\begin{eqnarray}}
\newcommand{\eeqn}{\end{eqnarray}}
\newcommand{\cred}[1]{\textcolor{red}{#1}}
\newcommand{\cblue}[1]{\textcolor{blue}{#1}}
\newcommand{\yzy}[1]{{\color[rgb]{0.70,0.13,0.13}{\textsf{YZY}: #1}}}
\newcommand{\hwd}[1]{{\color[rgb]{0.20,0.18,0.8}{\textsf{HWD}: #1}}}

\begin{document}
\setlength{\parindent}{2em}
\renewcommand{\refname}{参考文献}
\renewcommand{\figurename}{图}
\renewcommand{\tablename}{表}
\renewcommand{\acknowledgmentsname}{致谢}

\title{基于可训练 token 嵌入的张量网络 Born 机器上的序列学习\\
（Sequential learning on a Tensor Network Born machine with Trainable Token Embedding）}

\author{Wanda Hou}
\affiliation{Department of Physics, University of California, San Diego, CA 92093, USA}
\author{Miao Li}
\affiliation{School of Physics, Zhejiang University, Hangzhou 310027, China}
\affiliation{Department of Physics, Georgia Institute of Technology,
Atlanta, GA 30332, USA}
\author{Yi-Zhuang You}
\thanks{Corresponding author: yzyou@physics.ucsd.edu}
\affiliation{Department of Physics, University of California, San Diego, CA 92093, USA}
\date{\today}

\begin{abstract}
 生成模型（generative model）旨在学习数据背后潜在的概率分布，从而生成新的、逼真的样本。量子启发式生成模型，例如基于矩阵乘积态（MPS）框架的 Born 机器（Born machine），已在无监督学习任务中展现出卓越的能力。本研究通过正算符取值测量（POVM）引入可训练的 token 嵌入，取代传统的静态张量指标方案，从而推进了 Born 机器范式。关键技术创新包括：将 token 编码为带可训练参数的量子测量算符，并利用 QR 分解调整 MPS 的物理维数（physical dimension）。这一方法最大化了算符空间的利用率，并增强了模型的表达能力。在 RNA 数据上的实证结果表明，与 one-hot 嵌入相比，所提方法显著降低了负对数似然（NLL）；更高的物理维数进一步提升了单点概率与多点关联。该模型在单点估计上优于 GPT-2，并取得了具有竞争力的关联建模效果，展示了可训练 POVM 嵌入在量子启发式序列建模中刻画复杂数据关联的潜力。
\end{abstract}

\maketitle

\section{引言}

无监督学习对于充分利用当今海量的无标注数据至关重要。一种重要的方法是生成建模（generative modeling），其目标是在假定某种潜在物理机制的前提下对数据的概率分布进行建模\cite{lecun2015deep,2019RvMP...91d5002C}。与旨在寻找不同类别之间决策边界的判别模型不同，生成模型的目标是复现观测数据潜在的分布 $p(x)$，进而生成在原始数据集语境下合理的新样本。

近年来，生成建模技术取得了快速发展，许多模型，如变分自编码器（VAE）\cite{2013arXiv1312.6114K}、生成对抗网络（GAN）\cite{2014arXiv1406.2661G,2018PhRvA..98a2324D}以及 Transformer\cite{2017arXiv170603762V}，已被应用于各个领域。从经典物理中汲取灵感，现有模型或多或少都采用了基于能量的建模思想，即如同统计物理中那样，将某些概率建模为遵循 Boltzmann 分布（Boltzmann distribution）的 $p(x)\propto e^{-E(x)}$\cite{ackley1985learning,salakhutdinov2015learning}。与之相对，量子物理采用波函数（wave function）模型，其中概率被建模为 $p(x)=|\Psi(x)|^2$，其中 $\Psi(x)$ 是一个量子波函数。这类量子生成模型被称为 Born 机器\cite{PhysRevX.8.031012, 2022PhRvA.106b2612K, 2021arXiv211205326M}。量子概率模型比其经典对应物更具表达力，因为它们能够建模违反 Bell 不等式的分布，而后者正是量子纠缠（entanglement）的特征\cite{trueblood2012quantum,haven2016quantum}。

作为一种生成建模算法，Born 机器具有许多特色优势，例如可处理的 log-likelihood 与边缘化 log-likelihood，这提供了更好的可解释性并支持异常检测；log-likelihood 梯度可被高效且精确地求值，使得无需采样的基于梯度的优化成为可能；同时支持自回归（autoregressive）采样与掩码采样。凭借这些特性，Born 机器在无监督学习任务的许多领域（如序列数据学习）中表现出色，并能超越基于经典物理的模型\cite{2018Entrp..20..583C,2018arXiv180404168L}。

一类特别重要的数据是一维序列数据，即按照特定序列排序的一组数据点，其排序通常由时间或空间顺序决定。Born 机器在序列建模方面发展成熟，这是因为序列的一维结构允许使用一种高效的张量网络表示——矩阵乘积态（MPS）——来对底层量子态进行建模\cite{orus2019tensor,2019PhRvX...9c1041L,2021arXiv210612974L}。Born 机器支持可处理的 log-likelihood，从而提供更好的可解释性与高效的直接采样，同时支持自回归采样与掩码采样。目前的研究通过利用量子态的纠缠特征并将其编码到 MPS 中以提升 Born 机器的表达能力，已取得巨大进展\cite{2018PhRvA..98a2324D,mitarai2018quantum,2018arXiv180404168L,gong2022born,2019arXiv190300556M}。

对离散型序列数据 $\vect{x}=(x_1,x_2,\cdots)$ 进行建模时，相应的 token $x_i\in \dsN$ 通常被直接视为 MPS 物理维数的张量指标，并被编码为对角基投影算符 $\ket{x_i}\bra{x_i}$。然而在本研究中，我们探索一种更一般的嵌入方法：将每个 token $x_i$ 编码为局域 Hilbert 空间中带可训练参数 $\gamma$ 的更一般的量子测量算符 $M_{\gamma}(x_i)$，而非对角基投影算符。概率被建模为 $p_{\theta,\gamma}(\vect{x})=\bra{\Psi_\theta}M_{\gamma}(\vect{x})\ket{\Psi_\theta}$，其中 $M_\gamma(\vect{x})=\bigotimes_i M_\gamma(x_i)$ 是对应于序列 $\vect{x}$ 的联合测量算符，而 $\ket{\Psi_\theta}$ 是一个多体量子态，仍具有以 $\theta$ 参数化的 MPS 结构。关键创新在于：将 token 编码为测量算符充分利用了算符空间，使得在小得多的 Hilbert 空间中能够容纳更多 token。此外，通过同时提高嵌入量子测量与 MPS 的物理维数，模型可以利用更大算符空间的力量来提升表达能力。我们在序列 RNA 数据上的实验结果表明，借助调整物理维数的灵活性，Born 机器能够展现出更好的性能，并从数据集中学到更深层的相关性。

\section{编码结构（Encoding structure）}

生成建模的目标是在给定训练数据集的条件下对联合概率 $p(\vect{x})$ 进行建模。一旦模型训练良好，它就可以用来生成理想情况下应与训练数据分布相似的新样本。在本工作中，我们采用传统的 Born 机器（Born machine）结构，它以量子波函数 $\ket{\Psi_\theta}$ 为 backbone，并将编码 token 作为量子测量算符 $M(\vect{x})=\{M_{\gamma}(x_i)|x_i\in V\}$，其中 $\gamma$、$\theta$ 分别是来自嵌入（embedding）与 MPS 的可训练参数，$V$ 是全部词表的集合。每个数据点 $\vect{x}=(x_1,x_2,...,x_n)$ 的联合概率可以通过量子测量 $M_\gamma(\vect{x})=\bigotimes_{i}M_{\gamma}(x_i)$ 在给定量子波函数下的期望值来嵌入：
\begin{equation}\label{eq: P}
    p_{\theta, \gamma}(\vect{x})=\bra{\Psi_\theta}\bigotimes_{i}M_{\gamma}(x_i)\ket{\Psi_\theta},
\end{equation}
其中 backbone 量子态是归一化的，即 $\braket{\Psi_\theta}{\Psi_\theta}=1$。为了保证编码后的概率相对于整个输入空间是\textit{实的}、\textit{半正定}且\textit{归一化}的，我们把嵌入算符限制为正算符取值测量（positive operator-valued measurements, POVM），它们是\textit{厄米}的、\textit{半正定}的，且总体上\textit{归一化}：
\begin{equation}\label{eq: emb}
\begin{split}
    \forall x_{i}\in V,\ M_{\gamma}(x_i)&=M^{\dagger}_{\gamma}(x_i)\succeq 0,\\
    \sum_{i}M_{\gamma}(x_i)&=I.
\end{split}
\end{equation}
其中 $I$ 表示恒等算符。先前的研究已经探索了自适应 POVM 的多种实现方式 \cite{2022arXiv220807817G, 2023arXiv230315353Z, 2021PRXQ....2d0342G}；本工作通过在 Born 机器的序列数据建模中对可训练 token 嵌入采用 QR 分解，将这一概念加以扩展。


在先前的 Born 机器模型参数化方案中，嵌入约束是通过把每个 token $x_i\in\dsN$ 直接指定为由 token 索引所标记的基态的相应投影算符 $\ket{x_i}\bra{x_i}$ 来满足的\cite{PhysRevX.8.031012}。这等价于把 POVM 固定为 one-hot 嵌入矩阵 $M_\text{one-hot}(x_i)_{\alpha\beta}=\delta_{\alpha,x_i}\delta_{\beta,x_i}$，其矩阵维数严格等于词表大小（vocabulary size）。虽然 one-hot 嵌入可以满足 \eqnref{eq: emb}，但这一嵌入方案是预先定义好的、不可训练的。它无法充分利用算符空间，因而是次优的。相比之下，可训练的 POVM 嵌入使得有可能在同一个 Hilbert 空间中容纳更多 token，并进一步把量子态的物理维数（physical dimension）作为额外的超参数来提升模型的表达能力。在本文中，我们提出一种可训练的参数化方法，它在满足 \eqnref{eq: emb} 的同时能够发挥算符空间的潜力。

参数 $\gamma$ 被初始化为形状为 $(v \times p, p)$ 的复张量，其中 $v$ 表示词表大小，$p$ 表示量子态的物理维数。随后对 $\gamma$ 施加 QR 分解 $\gamma=Q_\gamma R_\gamma$，并将得到的 $Q_{\gamma}$ 矩阵进一步重塑为 $(v, p, p)$，再与其厄米共轭做批量相乘，得到量子测量矩阵。
\begin{equation}\label{eq: parametrization}
    M_{\gamma}(x_{i})=Q^{\dagger}_{\gamma}(x_i)Q_{\gamma}(x_i), \ x_i\in V.
\end{equation}
上述构造保证了 \eqnref{eq: emb} 中的约束自动得到满足，并且物理维数成为一个可调节的超参数。

遵循 \cite{PhysRevX.8.031012}，backbone 量子态可以用归一化的 MPS（matrix product state，矩阵乘积态）高效参数化，并且在基于梯度的优化过程中，相应的等距（isometric）约束 $\braket{\Psi_\theta}{\Psi_\theta}=1$ 通过等距优化方法加以保持\cite{wen2013feasible,2022MLS&T...3a5020G}。在物理 Hilbert 空间中显式选取基 $\ket{\vect{\beta}}=\ket{\beta_1,\beta_2,\cdots,\beta_n}$ 后，MPS 波函数可以表示为
\begin{equation}\label{eq: mps}
\braket{\vect{\beta}}{\Psi_\theta}=\sum_{b_1 b_2...b_{n-1}}(A_{\theta 1})^{\beta_1}_{b_1} (A_{\theta 2})^{\beta_2}_{b_1 b_2}\dots (A_{\theta n})^{\beta_n}_{b_{n-1}},
\end{equation}
它由一系列在位点 $i$ 处以 $\theta$ 参数化的张量 $A_{\theta i}$ 来描述。张量 $A_{\theta i}$ 是等距（isometry）的，其含义是：在收缩 $A_{\theta i}$ 与 $A_{\theta i}^*$（复共轭）之间的入腿之后，出腿构成一个恒等映射，
\begin{equation}
\sum_{b_i,\beta_i} (A_{\theta i}^*)_{b_{i-1}b_{i}}^{\beta_i}(A_{\theta i})_{b'_{i-1}b_{i}}^{\beta_i}=\delta_{b_{i-1}b'_{i-1}.}
\end{equation}
等距结构在 \figref{fig: tn} 中由箭头标示。采用等距 MPS 有若干优点。首先，它保证了所建模概率分布的归一化，从而可以直接进行对数似然（log-likelihood）优化。其次，
它允许通过求和消去后续 token，把概率分布高效地边缘化（marginalization）到首个 token 上。

\begin{figure}[htbp]
\begin{center}
\includegraphics[width=240pt]{tn.pdf}
\caption{概率模型 $p_{\theta,\gamma}(\vect{x})$ 的张量网络表示。MPS 张量 $A_{\theta i}$ 为蓝色圆圈，测量算符 $M_\gamma(x_i)$ 为橙色方框。等距结构由张量腿上的箭头标示，表示从较大的 Hilbert 空间投影到较小的子空间。}
\label{fig: tn}
\end{center}
\end{figure}

给定 Born 机器易于处理的对数似然，针对指定数据集的无监督学习目标函数即为负对数似然（negative log-likelihood, NLL）\begin{equation}
\mathcal{L}=-\frac{1}{N}\sum_{\vect{x}\in \text{data}}\log p_{\theta,\gamma}(\vect{x}),
\end{equation}
其中 $N$ 表示数据集的大小。\figref{fig: tn} 展示了单个数据点按一维形状 $\vect{x}=(x_1,x_2,...,x_n)$ 进行概率编码的张量网络。基于 NLL 优化参数的训练过程在 算法~\ref{alg:train} 中给出。


\begin{algorithm}[H]
\caption{Born 机器的训练过程}
\label{alg:train}
\begin{algorithmic}
\REQUIRE 训练数据集 $D$，学习率（learning rate）$\eta$
\STATE 初始化参数 $\theta$，$\gamma$
\WHILE{未收敛}
    \FOR{每个数据样本 $x \in D$}
        \STATE 使用 \eqnref{eq: P} 计算联合概率 $p_{\theta,\gamma}(x)$
        \STATE 计算 NLL 损失 $L = -\log p_{\theta,\gamma}(x)$
        \STATE 计算梯度 $\nabla_{\theta} L$，$\nabla_{\gamma} L$
        \STATE 更新参数：
        \STATE $\theta \leftarrow \theta - \eta \nabla_{\theta} L$
        \STATE $\gamma \leftarrow \gamma - \eta \nabla_{\gamma} L$
    \ENDFOR
    \STATE 检查是否收敛
\ENDWHILE
\ENSURE 优化后的参数 $\theta$，$\gamma$
\end{algorithmic}
\end{algorithm}
\section{边缘化与生成式采样}

训练完成后，Born 机器可以按照 \eqnref{eq: P} 进行自回归生成；相比之下，energy-based model 通常需要借助 Markov 链蒙卡（MCMC）才能从 $p(\vect{x})=\frac{1}{Z}e^{-E(\vect{x})}$ 中采样。Born 机器的自回归采样过程逐位进行，依赖于其计算条件概率/边缘概率的能力。重要的是，由于张量收缩运算满足交换律，采样方向不受限制，既不必从头开始，也可以沿任意方向进行。因此，对于已给定 token 的位点，可以通过插入相应的 POVM 对其施加条件，然后再求迹移除；而对于待采样的位点，则可以直接对 MPS 相应的物理腿求迹移除来实现边缘化，这得益于 \eqnref{eq: emb} 中的归一化性质。整个序列的生成通过迭代地重复同一过程、填满所有空位来完成。

例如，设 $K\subseteq\{1,2,...,n\}$ 为一个位点子集，$\vect{x}_{K}=\{x_i\}_{i\in K}$ 为区域 $K$ 内的 token。训练完成后，假设我们想在给定补集区域 $\bar{K}$ 中其余 token $\vect{x}_{\bar{K}}$ 的条件下，计算条件分布 $p_{\theta,\gamma}(\vect{x}_K|\vect{x}_{\bar{K}})$。它可以按下式求值：
\begin{equation}\label{eq: single site}
    p_{\theta,\gamma}(\vect{x}_K|\vect{x}_{\bar{K}})=\bra{\Psi_\theta(\vect{x}_{\bar{K}})}M_{\gamma}(\vect{x}_K)\ket{\Psi_\theta(\vect{x}_{\bar{K}})},
\end{equation}
其中 $M_{\gamma}(\vect{x}_K)=\bigotimes_{i\in K}M_{\gamma}(x_i)$ 是区域 $K$ 内的测量算符，而 $\ket{\Psi_\theta(\vect{x}_{\bar{K}})}$ 表示根据区域 $\bar{K}$ 中的测量结果 $\vect{x}_{\bar{K}}$ 坍缩后的量子态，其定义为
\begin{equation}
\ket{\Psi_\theta(\vect{x}_{\bar{K}})}=\frac{1}{Z}\bigotimes_{i\in\bar{K}}Q_{\gamma}(x_i)\ket{\Psi_{\theta}},
\end{equation}
其中 $Z=\bra{\Psi_\theta}M(\vect{x}_{\bar{K}})\ket{\Psi_\theta}$ 是该态的归一化常数。
与常规自回归模型只能沿因果顺序计算条件概率分布不同，这里的 Born 机器能够计算任意子集 $K$ 中 token 的概率分布，其条件为补集 $\bar{K}$ 中的 token。
% Based on this, the stand-by site has the single site probability of each token
% \begin{equation}\label{eq: single site}
%     p_{\theta,\gamma}(x_i|\vect{x}_0)=\bra{\Psi_\theta(\vect{x}_0)}M_{\gamma}(x_i)\ket{\Psi_\theta(\vect{x}_0)},
% \end{equation}
% where $\ket{\Psi(\vect{x}_0)}$ is obtained by acting the conditional POVMs onto the backbone quantum state with arbitrary order. \yzy{What is the mathematical definition of $\Psi(x_0)$? Unlike autoregressive models, which are greedy, token sampling must be performed following a causal order. Born machine, each measurement collapses the wave function globally, and token sampling can be performed in any order.} \hwd{I think $\Psi(x_0)$ is given by collapsing the original wave function according to measurement outcomes, $\ket{\Psi_\theta(\vect{x}_0)}=\mathcal{P}(\vect{x}_0)\ket{\Psi_\theta}=\mathcal{P}(x_{v_1,i_1}, x_{v_2,i_2},\dots, x_{v_m,i_m})\ket{\Psi_\theta}$ where $v, i$ label the vocabulary index/site index, and $m$ is the number of measured sites. Since each of the measure collapse the wave function globally, the projection procedure can be performed in arbitrary order for each site. In our case, the conditional environment of a projected site is calculated by inserting a POVM between the wave functions and the conjugate of itself, therefore the explicit form of our projection operator should be the matrix square-root of POVMs $\mathcal{P}(x_{v_1,i_1}, x_{v_2,i_2},\dots, x_{v_m,i_m})=\sqrt{M(x_{v_1,i_1})}\sqrt{M(x_{v_2,i_2})}\cdots\sqrt{M(x_{v_m,i_m})}$. For masked sampling task\eqref{eq: single site}, $\ket{\Psi_\theta(\vect{x}_0)}$ is further normalized; for joint probability evaluation procedure\eqref{eq: P}, the half-collapsed wave function in the intermediate is not normalized because the result value represents model $p_{\theta}(\vect{x}) \in (0,1)$.}

\begin{algorithm}[H]
\caption{带条件依赖的生成式采样}
\label{alg:gen}
\begin{algorithmic}
\REQUIRE 训练好的参数 $\theta$、$\gamma$，序列长度 $n$，可选的初始序列 $x_{\text{init}}$
\IF{未提供 $x_{\text{init}}$}
    \STATE 初始化长度为 $n$ 的序列 $x = [ \text{empty} ]$
\ELSE
    \STATE 以 $x_{\text{init}}$ 作为起始序列，其中元素可以为空或预先给定的 token
\ENDIF
\STATE 定义集合 $K$ 为已知 token 的下标集合，$\bar{K}$ 为待采样 token 的下标集合
\WHILE{$\bar{K}$ 非空}
    \FOR{$\bar{K}$ 中的每个下标 $i$}
        \IF{$x[i]$ 为空}
            \STATE 利用 \eqnref{eq: single site}，在给定已知 token 的条件下计算 token $x_i$ 的条件概率：
            \STATE $p(x_i|x_{K}) = \langle \Psi_{\theta}(x_{K}) | M_{\gamma}(x_i) | \Psi_{\theta}(x_{K}) \rangle$
            \STATE 从 $p(x_i|x_{K})$ 中采样得到 $x_i$
            \STATE 用采样得到的 token $x_i$ 更新 $x[i]$
            \STATE 将 $i$ 加入集合 $K$，并从 $\bar{K}$ 中移除 $i$
        \ENDIF
    \ENDFOR
\ENDWHILE
\ENSURE 生成的序列 $x$
\end{algorithmic}
\end{algorithm}

由于 Born 机器支持按任意顺序进行自回归生成，可以利用这一特性，通过计算单位点概率 \eqnref{eq: single site} 或多位点的联合概率，直接进行掩码采样或异常检测。此外，还可以将模型学到的边缘概率与训练数据集的统计频率进行比较，以评估训练后模型的性能，这将在 \secref{Empirical result} 中进一步讨论。该过程的详细实现见算法~\ref{alg:gen}。
\section{实验结果}\label{Empirical result}

由于 MPS 的一维结构，Born 机器（Born machine）特别适合于序列建模。在各种类型的序列数据中，生物基因序列是最为重要的一类，对人类具有巨大价值。然而，基因序列难以通过实验获得，且为提取有限的数据往往需要付出巨大成本。因此，使用生成模型来辨识基因序列中潜在的概率分布并生成新样本，可以提高探索新基因序列的效率并降低成本。在本研究中，我们将评估带可训练 POVM 的 Born 机器在 RNA 序列无监督学习中的性能。我们使用来自开源数据库 \textit{5SrRNAdb}\cite{szymanski20165srnadb} 的细菌 5S rRNA 数据，其平均长度为 120，词表总大小为 4（\textit{A}: 腺嘌呤 Adenine，\textit{T}: 胸腺嘧啶 Thymine，\textit{C}: 胞嘧啶 Cytosine，\textit{G}: 鸟嘌呤 Guanine）。训练集和测试集是从该数据集中无偏划分得到的。

\begin{figure}[htbp]
\begin{center}
\includegraphics[width=250pt]{samples.pdf}
\caption{在此图中，每条水平线代表一条 RNA 序列，纵向排列显示。四种颜色代表四种类型的核苷酸。(a) 由带 POVM 嵌入的 Born 机器生成的样本。(b) 由 GPT-2 生成的样本。(c) 来自测试集的真实数据。}
\label{fig: generation}
\end{center}
\end{figure}

\figref{fig: generation} 中由带 POVM 的 Born 机器与 GPT-2 生成的样本在定性上与真实测试集相似，表明两个模型都捕捉到了序列数据中可见的模式。鉴于这种相似性，我们接下来使用诸如 NLL、单位点概率（single-site probability）与 Pearson 相关（Pearson correlation）等定量指标来评估各模型的性能。\figref{fig: loss} 表明，增大 POVM 的物理维数（physical dimension）可以降低收敛的 NLL 值，即使在边界维数固定的情况下也是如此。此外，它还揭示出更高的物理维数扩展了键维（bond dimension）的上限，从而能够进一步最小化 NLL。

\begin{figure}[htbp]
\begin{center}
\includegraphics[width=200pt]{loss.pdf}
\caption{使用 \textit{Adam} 优化器最小化 NLL 目标函数，并为键维/物理维数分别调节学习率。结果表明，同时增大键维和物理维数可以降低收敛的 NLL。绿色虚线表示参数量相近的 GPT-2 模型的收敛 NLL。}
\label{fig: loss}
\end{center}
\end{figure}

RNA 数据分布的特征之一是，对各个位点做边缘化后的联合概率分布蕴含着潜在的接触预测（contact prediction）信息\cite{trinquier2021efficient}。RNA 中的接触预测是指确定 RNA 分子中哪些残基（residue）或核苷酸（nucleotide）在分子的三维结构中彼此邻近的过程。为了更深入地评估带可训练 POVM 的 Born 机器的性能并发掘接触预测，可以利用其高效的边缘化能力 \eqnref{eq: single site}，然后将其与从数据集中提取的统计频率进行比较。特别地，我们将研究两个\textit{统计特征}（statistical features）：
\begin{itemize}
\item 单位点概率分布 $p(x_i)$，
\item 以及两位点 Pearson 相关
\begin{equation}\label{eq: corre}
    c(x_i,x_j)=\log\frac{p(x_i,x_j)}{p(x_i)p(x_j)}.
\end{equation}
\end{itemize}
\figref{fig: corre} 中的每个子图画出模型预测的统计特征（纵轴）相对于数据给出的统计特征（横轴）的散点。\figref{fig: corre}(a-e) 对应 $p(x_i)$，其中图中每个点对应位点 $i$ 及该位点上 token $x_i$ 的不同选择。\figref{fig: corre}(f-j) 对应 $c(x_i,x_j)$，其中图中每个点对应配对 $(i,j)$ 及该配对上 token $(x_i,x_j)$ 的选择。如果模型完美地学到了数据中的所有统计特征，我们应当期望所有点都落在对角虚线上（在该线上模型预测与数据统计相一致）。\figref{fig: corre} 中的结果表明，随着物理维数的增大，单位点概率与关联都更加接近测试集。这说明，与固定的 one-hot 嵌入方法相比，配备可训练 POVM 嵌入的 Born 机器在建模训练数据中跨位点的 token 关联方面更为强大。


我们还将我们的量子模型与参数量相当的经典序列生成模型 Generative Pretrained Transformer-2（GPT-2）进行了基准比较。GPT-2 配置由 \textit{Hugging Face transformer} 库\cite{gpt2lmhead} 初始化，取 $d_{\text{embedding}} = 72$、$n_{\text{layer}} = 12$、$n_{\text{head}} = 4$，总共约含 80 万参数，并与 GPT-2LMHead 模型集成。对于 Born 机器，参数总数可以估计为键维、物理维数与序列长度的乘积，即 $d_{\text{bond}}^2 \times d_{\text{phys}} \times L$。例如，我们最大的配置使用 $d_{\text{bond}} = 20$、$d_{\text{phys}} = 16$、$L = 120$，得到约 77 万参数——与 GPT-2 相当。在 \figref{fig: loss} 与 \figref{fig: corre} 中，在我们当前的设置下，Born 机器相比 GPT-2 达到了更低的 NLL，并在近似单位点概率方面表现相当；但它在关联近似方面的能力较弱。像 GPT-2 这类基于注意力的网络具有全连接（all-to-all connectivity）结构，且没有内在的维数限制，因而特别擅长捕捉长程关联（long-range correlation）。

\begin{figure}[htbp]
\begin{center}
\includegraphics[width=240pt]{corre.pdf}
\caption{比较模型的单位点概率 $p(x_i)$（上排，(a-e)）与两位点关联 $c(x_i,y_j)$（下排，(f-j)）与数据集中对应的量。GPT-2：(a)\&(f)。one-hot 嵌入（基线）：(b)\&(g)。可训练 POVM 嵌入（本文方法），物理维数：(c)\&(h) $p=4$，(d)\&(i) $p=8$，(e)\&(j) $p=16$。虚线表示 $y=x$ 参考线。}
\label{fig: corre}
\end{center}
\end{figure}


\section{总结}

生成模型是机器学习的一个重要分支，它致力于理解数据固有的概率分布。与判别模型不同，它们侧重于捕捉数据的内在结构。本工作继续研究一类称为 Born 机器的量子启发式生成建模模型。这些机器在矩阵乘积态（MPS）框架内使用量子波函数对潜在的数据分布进行建模。本研究的主要创新在于对 token 嵌入方法的扩展。这项工作没有通过相应的张量指标直接嵌入每个 token，而是引入了可训练的正算符取值测量（POVM）的概念。通过将 MPS 与可训练嵌入相结合，Born 机器可以获得更好的性能，并从数据集中提取更深层的关联。源代码与原始数据可在 GitHub 仓库\cite{hou2023github} 中获取。

我们的新 Born 机器架构以量子波函数为主干，将 token 编码为量子测量矩阵。每个数据点的联合概率表示为一个量子测量的期望值。Born 机器的一个关键特性是其自回归生成能力，可以按任意顺序进行序列采样——这使其适用于掩码 token 采样（masked token sampling）与异常检测（anomaly detection）等任务。虽然我们当前的实现可以通过 MPS 进行经典模拟，但真正的优势可能要等到在真实量子设备上进行采样时才会显现。一个有前景的策略是先在经典硬件上用 MPS 高效地训练模型，然后通过等价的量子线路（quantum circuit）将训练好的态转移到量子计算机上。如 \cite{2022PhRvX..12a1047H, 2024QS&T....9a5012R} 所示范的，量子线路中的每个 MPS 张量都可以映射为带有 ancilla 插入与测量的局域量子门，从而将 MPS 表示为量子计算机上的真实量子态。此外，本工作引入的自适应 POVM 方法可以通过扩展 Hilbert 空间上的广义测量来实现\cite{preskill1998lecture}，从而实现高效采样并潜在地利用量子优势（quantum advantage）。由于 POVM 在实践中可以直接实现，这一方法可以扩展到 Quantum Circuit Born Machine (QCBM) 架构，以便在量子硬件上实现\cite{PhysRevA.98.062324, 2019npjQI...5...45B}。

在应用方面，我们的研究聚焦于 RNA 序列的无监督学习。结果表明，通过增大 POVM 中的物理维数，Born 机器可以获得更好的性能，捕捉 RNA 数据中更深层的关联。然而，Born 机器的当前框架中仍存在一些未解决的问题。例如，其现有结构无法直接处理长度差异显著的数据。此外，由于当前的 Born 机器是用一维 MPS 与离散化的嵌入方法构建的，直接将其扩展到处理更高维与连续的数据类型并不方便。另外，还存在将训练好的 MPS 与相应的 POVM 转移到量子设备上进行高效采样的挑战。

\begin{acknowledgments}
我们感谢与 Rose Yu、Tian-Xiao Hu 和 Zhao-Yi Zeng 的有益讨论。
我们感谢 OpenAI GPT4 模型在撰写本文的过程中提供编辑建议。本研究由 NSF Grant No.~DMR-2238360 资助。
\end{acknowledgments}


\bibliographystyle{apsrev4-2}
\bibliography{ref}












\onecolumngrid
\newpage
\appendix
\setcounter{secnumdepth}{2}


\end{document}

